\documentclass[../../main.tex]{subfiles}
\begin{document}
\section{同伦与基本群}\label{section:appA-homotopy-fundamental-group}

如果 $X$ 与 $Y$ 是拓扑空间，$F_0,F_1:X\to Y$ 是连续映射，则从 $F_0$ 到 $F_1$ 的一个\emph{同伦}是连续映射 $H:X\times I\to Y$，满足对所有 $x\in X$ 都有
\begin{align*}
H(x,0)=F_0(x),\qquad H(x,1)=F_1(x).
\end{align*}
如果存在从 $F_0$ 到 $F_1$ 的同伦，就说 $F_0$ 与 $F_1$ \emph{同伦}，写作 $F_0\simeq F_1$。如果同伦满足对所有 $t\in I$ 和某子集 $A\subset X$ 中所有 $x$ 都有 $H(x,t)=F_0(x)=F_1(x)$，则称映射 $F_0$ 与 $F_1$ \emph{相对于 $A$ 同伦}。“同伦”和“相对于 $A$ 同伦”都是从 $X$ 到 $Y$ 的所有连续映射集合上的等价关系。

同伦最重要的应用是道路。假设 $X$ 是拓扑空间。两条道路 $f_0,f_1:I\to X$ 如果相对于 $\{0,1\}$ 同伦，就称它们\emph{道路同伦}，符号记作 $f_0\sim f_1$。明确地说，这意味着存在连续映射 $H:I\times I\to X$ 满足
\begin{align*}
H(s,0)=f_0(s),\quad s\in I,\qquad
H(s,1)=f_1(s),\quad s\in I,
\end{align*}
\begin{align*}
H(0,t)=f_0(0)=f_1(0),\quad t\in I,\qquad
H(1,t)=f_0(1)=f_1(1),\quad t\in I.
\end{align*}
对任意给定的点 $p,q\in X$，道路同伦是从 $p$ 到 $q$ 的所有道路集合上的等价关系。道路 $f$ 的等价类称为它的\emph{道路类}，记作 $[f]$。

给定两条道路 $f,g:I\to X$ 满足 $f(1)=g(0)$，它们的\emph{积}是道路 $f\cdot g:I\to X$，定义为
\begin{align*}
(f\cdot g)(s)=
\begin{cases}
f(2s),&0\leqslant s\leqslant\tfrac12,\\
g(2s-1),&\tfrac12\leqslant s\leqslant1.
\end{cases}
\end{align*}
如果 $f\sim f'$ 且 $g\sim g'$，不难证明 $f\cdot g\sim f'\cdot g'$。因此定义道路类 $[f]$ 与 $[g]$ 的积 $[f]\cdot[g]=[f\cdot g]$ 是有意义的。虽然道路乘法不结合，但它在道路同伦意义下是结合的：$([f]\cdot[g])\cdot[h]=[f]\cdot([g]\cdot[h])$。当我们需要考虑三条或更多条真实道路（而非道路类）的积时，我们约定这样的积从左到右计算：$f\cdot g\cdot h=(f\cdot g)\cdot h$。

如果 $X$ 是拓扑空间而 $q$ 是 $X$ 中一点，$X$ 中\emph{以 $q$ 为基点的闭路}是 $X$ 中从 $q$ 到 $q$ 的道路，即连续映射 $f:I\to X$ 满足 $f(0)=f(1)=q$。以 $q$ 为基点的闭路的道路类集合记作 $\pi_1(X,q)$。配以上述乘法，它是一个群，称为 $X$ \emph{以 $q$ 为基点的基本群}。该群的单位元是常值道路 $c_q(s)\equiv q$ 的道路类，$[f]$ 的逆是反向道路 $\bar f(s)=f(1-s)$ 的道路类。

可以证明，对道路连通空间，以不同点为基点的基本群同构。如果 $X$ 道路连通，且对某个（因而每个）$q\in X$，基本群 $\pi_1(X,q)$ 是仅由 $[c_q]$ 组成的平凡群，就说 $X$ \emph{单连通}。这意味着每条闭路都道路同伦于一条常值道路。

\begin{exercise}\label{exercise:appA-A-61}
设 $X$ 是道路连通拓扑空间。证明 $X$ 单连通，当且仅当 $X$ 中每对起点和终点相同的道路都道路同伦。
\end{exercise}

同伦关系的一个关键特征是它在复合下保持，如下一个命题所示。

\begin{proposition}\label{proposition:appA-homotopy-composition}
如果 $F_0,F_1:X\to Y$ 与 $G_0,G_1:Y\to Z$ 是连续映射，且 $F_0\simeq F_1$、$G_0\simeq G_1$，则 $G_0\circ F_0\simeq G_1\circ F_1$。类似地，如果 $f_0,f_1:I\to X$ 道路同伦而 $F:X\to Y$ 连续，则 $F\circ f_0\sim F\circ f_1$。
\end{proposition}

\begin{exercise}\label{exercise:appA-A-63}
证明前一个命题。
\end{exercise}

因此，如果 $F:X\to Y$ 是连续映射，对每个 $q\in X$，我们得到一个良定义映射
\begin{align*}
F_*:\pi_1(X,q)\to\pi_1(Y,F(q))
\end{align*}
由 $F_*[f]=[F\circ f]$ 给出。

\begin{proposition}\label{proposition:appA-induced-hom}
如果 $X$ 与 $Y$ 是拓扑空间而 $F:X\to Y$ 连续，则 $F_*:\pi_1(X,q)\to\pi_1(Y,F(q))$ 是群同态，称为由 $F$ \emph{诱导的同态}。
\end{proposition}

\begin{proposition}[诱导同态的性质]\label{proposition:appA-induced-hom-properties}
\begin{enumerate}[(a)]
\item 设 $F:X\to Y$ 与 $G:Y\to Z$ 是连续映射。则对每个 $q\in X$，$(G\circ F)_*=G_*\circ F_*:\pi_1(X,q)\to\pi_1(Z,G(F(q)))$。
\item 对每个空间 $X$ 和每个 $q\in X$，恒等映射 $\operatorname{Id}_X:X\to X$ 诱导的同态是 $\pi_1(X,q)$ 的恒等映射。
\item 如果 $F:X\to Y$ 是同胚，则 $F_*:\pi_1(X,q)\to\pi_1(Y,F(q))$ 是同构。因此，同胚空间有同构的基本群。
\end{enumerate}
\end{proposition}

\begin{exercise}\label{exercise:appA-A-66}
证明前两个命题。
\end{exercise}

\begin{exercise}\label{exercise:appA-A-67}
如果 $U\subset\mathbb R^n$ 有一点 $c\in U$，使对每个 $x\in U$，从 $c$ 到 $x$ 的线段都含于 $U$，就称 $U$ \emph{星形凸}。证明每个星形凸集单连通。
\end{exercise}

\begin{proposition}[球面的基本群]\label{proposition:appA-sphere-fundamental}
\begin{enumerate}[(a)]
\item $\pi_1(S^1,(1,0))$ 是由闭路 $\omega:I\to S^1$（给定为 $\omega(s)=(\cos2\pi s,\sin2\pi s)$）的道路类生成的无限循环群。
\item 如果 $n>1$，$S^n$ 单连通。
\end{enumerate}
\end{proposition}

\begin{proposition}[积空间的基本群]\label{proposition:appA-product-fundamental}
假设 $X_1,\dots,X_k$ 是拓扑空间，$p_i:X_1\times\cdots\times X_k\to X_i$ 表示第 $i$ 个投影映射。对任意点 $q_i\in X_i$（$i=1,\dots,k$），定义映射
\begin{align*}
P:\pi_1(X_1\times\cdots\times X_k,(q_1,\dots,q_k))
\to\pi_1(X_1,q_1)\times\cdots\times\pi_1(X_k,q_k)
\end{align*}
由
\begin{align*}
P[f]=(p_{1*}[f],\dots,p_{k*}[f])
\end{align*}
给出。则 $P$ 是同构。
\end{proposition}

\begin{exercise}\label{exercise:appA-A-70}
证明前两个命题。
\end{exercise}

拓扑空间之间的连续映射 $F:X\to Y$ 如果存在连续映射 $G:Y\to X$ 使 $F\circ G\simeq\operatorname{Id}_Y$ 且 $G\circ F\simeq\operatorname{Id}_X$，就称为\emph{同伦等价}。这样的映射 $G$ 称为 $F$ 的一个\emph{同伦逆}。如果 $X$ 与 $Y$ 之间存在同伦等价，就说这两个空间\emph{同伦等价}。例如，包含映射 $\iota:S^{n-1}\hookrightarrow\mathbb R^n\setminus\{0\}$ 是一个同伦等价，同伦逆为 $r(x)=x/|x|$，因为 $r\circ\iota=\operatorname{Id}_{S^{n-1}}$，而 $\iota\circ r$ 通过直线同伦 $H(x,t)=tx+(1-t)x/|x|$ 同伦于 $\mathbb R^n\setminus\{0\}$ 的恒等映射。

\begin{theorem}[同伦不变性]\label{theorem:appA-homotopy-invariance}
如果 $F:X\to Y$ 是同伦等价，则对每个 $p\in X$，$F_*:\pi_1(X,p)\to\pi_1(Y,F(p))$ 是同构。
\end{theorem}

证明见本附录开头提到的任一拓扑学教材。
\end{document}
